\section{Naming Conventions}
\label{sec:NamingConventions}

Durga relies heavily on predictable names because the generated topology is intentionally
visible. Naming is therefore part of the system design, not only a cosmetic detail.

\subsection{Common naming patterns}

The generator generally follows these patterns:
\begin{itemize}
\item process identifiers come from BPMN process ids,
\item task topics use process and task identifiers,
\item generated classes use process and BPMN element names converted to Java class names,
\item helper scripts use short imperative file names such as \texttt{complete-task.sh},
\item the canonical lifecycle-event topic is per process by default,
\filepath{process-events-\%processId\%} (override with \texttt{--event-topic});
the monitor consumes the whole \filepath{process-events-*} family,
\item shared monitoring output topics use fixed names such as:
\begin{quote}
\filepath{process-state}\\
\filepath{process-state-counts}
\end{quote}
\item validation mode uses a dedicated consumer group
\texttt{\%processId\%\_\%taskId\%\_validation} on the production input topic and the
shared topics \filepath{validation-candidate-outputs} and \\ \filepath{validation-results}
(Section~\ref{sec:ValidationMode}).
\end{itemize}

\subsection{Why this matters}

Consistent names make it easier to:
\begin{itemize}
\item inspect topic traffic in Kafka tools,
\item map BPMN elements to generated Java handlers,
\item reason about subprocess and boundary scope behavior,
\item debug routing and completion behavior quickly.
\end{itemize}

\subsection{Practical guidance}

For the cleanest generated output, BPMN ids should be stable, readable and
explicit. If BPMN models use unclear or heavily tool-generated ids, the
generated Kafka and Java output becomes harder to read as well.

If the BPMN tool changes element identifiers between edits, pin the process
id with \texttt{--process-id <id>} to keep topic names, consumer groups,
and class names stable across regenerations.
